% ============================ APPENDICES ============================

\section{Notation reference}\label{app:notation}

One symbol denotes one concept throughout. Entries marked \textbf{[clash]} are
symbols that collide across the source literature and are disambiguated here.

\renewcommand{\arraystretch}{1.15}
\begin{longtable}{@{}p{0.42\textwidth} p{0.52\textwidth}@{}}
\toprule
\textbf{Symbol} & \textbf{Concept} \\
\midrule
\endhead
\multicolumn{2}{@{}l}{\textit{Model, kernel, structured prior}}\\
\midrule
$x$, $x^{*}$ & stimulus image (vectorized natural image); candidate image \\
$\pixc_i$, $\rfc$ & pixel coordinate; RF center $(\varepsilon_{0x},\varepsilon_{0y})$ \\
$\lat(x)$ & latent GP function (the tuning latent) \textbf{[clash: not $\latz$]} \\
$\latz$ & firing-rate offset \\
$\rate(x)=e^{\gain\lat+\latz}$ & firing rate \textbf{[clash: $\rate$ is the rate, never the latent]} \\
$\lgr=\gain\lat+\latz$ & log-firing rate (so $\rate=e^{\lgr}$) \\
$\gain$ & Poisson gain \textbf{[clash: not the projection $\amat$]} \\
$\spk$ ($\equiv r$) & spike count (synonyms $r,n$, \texttt{target}) \\
$\Kf(x,x')$ & arc-cosine prior kernel function \\
$\Ktil=\Kf(\Ipt,\Ipt)$ & inducing gram \textbf{[clash: never a bare $K$ for this]} \\
$\Kmag=\sqrt{\vmag_x\vmag_{x'}}$ & kernel magnitude \textbf{[clash: not $\Mind$]} \\
$\kang$, $\Jang(\kang)$ & kernel angle; angular term \textbf{[$\Jang$ is never a Jacobian]} \\
$\vmag_x=x\transpose\Cmat x+\sigz^{2}$ & kernel self-magnitude \\
$\Cmat$, $\Csm$ & structured spatial prior; RBF smoothing factor \\
$\loc_i$ & locality mask / RF envelope \textbf{[clash: not diffusion $\aret$]} \\
$\bwid$ & RF half-width \textbf{[clash: not diffusion $\bsched$]} \\
$\smoo$ & smoothness SD \textbf{[clash: not correlation $\corr$]} \\
$\sigz$ & bias std, $\sigz=e^{\mathrm{raw\_sigma\_0}}$ \\
$\Amp$ & amplitude --- kernel hyperparameter, $\Amp=\mathrm{softplus}(\text{raw\_Amp})$, bounds $(0,1000]$, optimized by default \\
\midrule
\multicolumn{2}{@{}l}{\textit{Variational inference and eigenspace}}\\
\midrule
$\ipt$, $\Ipt$ & inducing point; inducing-point matrix \\
$\latt$ ($\equiv u$) & inducing latent values, $q(\latt)=\Gauss(\vm,\vV)$ \\
$\Mind$, $\Ntr$ & number of inducing points; number of training images ($\Mind{=}\Ntr$ in the loop) \\
$\kvec(x)$, $\pv(x)=\Ktil^{-1}\kvec(x)$ & cross-covariance vector; projection vector \\
$\vm$, $\vV$ & variational mean and covariance (inducing space) \\
$\pmean(x)$, $\pvar(x)$ & posterior mean and variance of the latent \\
$\Lchol$ & Cholesky factor of $\Ktil$ \textbf{[clash: $\ELBO$ is the ELBO]} \\
$\eb$, $\evals$, $\nb$ & eigenbasis of $\Ktil$; eigenvalues (diagonal); kept dimension \textbf{[$\evals$ not $\lat$]} \\
$\mb$, $\Vb$, $\Ktilb$ & eigenspace mean, dense covariance, diagonal inducing gram \\
$\amat=\Kf\Ktil^{-1}$ & projection matrix (\texttt{KKtilde\_inv\_b}) \\
$\fbar=e^{\gain\pmean+\frac12\gain^{2}\pvar+\latz}$ & expected firing rate (\texttt{f\_mean}) \\
$\ELBO$ & evidence lower bound \textbf{[clash: script $\mathcal L$, not the Cholesky factor]} \\
\midrule
\multicolumn{2}{@{}l}{\textit{Training-step helpers}}\\
\midrule
$\gE=\gain\sum_i\kvec_i(\spk_i-\fbar_i)$ & E-step gradient helper \textbf{[clash: not the log-rate $\lgr$]} \\
$\GE=\gain^{2}\sum_i\fbar_i\kvec_i\kvec_i\transpose$ & E-step Fisher curvature (PSD) \\
\midrule
\multicolumn{2}{@{}l}{\textit{Utility and information}}\\
\midrule
$\Ustd$, $\Uda$ & standard (production) and distribution-aware (research) utility \\
$\Hmarg$, $\Hnoise$, $\Hcond$ & marginal, aleatoric-noise, and conditional entropy \textbf{[clash: never a bare $H$]} \\
$\px$ & natural-image distribution \\
$\mug=\gain\pmean+\latz$, $\varg=\gain^{2}\pvar$ & log-rate mean and variance \\
$\corr(x,x^{*})$ & posterior latent correlation \textbf{[clash: not smoothness $\smoo$; $\corr^{2}$ = fractional variance reduction]} \\
$\ccond$ & prior conditional covariance (the ``correction'' term) \\
$\scal$ & input scaling factor, $x^{*}=\scal\,x_t$ \textbf{[clash: not $\ccond$]} \\
$\varr$, $\subc$ & residual/aleatoric variance; subspace coordinate \textbf{[$\subc$ not an inducing point]} \\
$\LW$, $\gmode$, $\rmax$ & Lambert-$W$ (principal); Laplace mode; count truncation \\
$\satf$ & saturation-kernel amplitude \\
\midrule
\multicolumn{2}{@{}l}{\textit{Diffusion (Part~IV only)}}\\
\midrule
$\dt$ & diffusion timestep \textbf{[there is no temporal RF: spatial GP only]} \\
$\xz$, $\xt$, $\xhat$ & clean, noisy, and Tweedie-estimated image \\
$\bsched$ & noise-schedule variance \textbf{[clash: not GP $\bwid$]} \\
$\aret$, $\abar$ & signal retention; cumulative retention \textbf{[clash: not GP mask $\loc$]} \\
$\epsth$, $\score$ & U-Net noise predictor; score function \\
$\gw$, $\ggrad$ & guidance scale; guidance gradient ($=$ the GP utility gradient) \\
$\greg$ & approach-A regularization weight \textbf{[clash: not the GP latent $\lat$]} \\
\bottomrule
\end{longtable}

\clearpage
\section{Reconciled discrepancies and open caveats}\label{app:discrepancies}

The document is anchored to the current code. The following notes-vs-code and
notes-vs-notes discrepancies were reconciled during consolidation (code wins),
and the honest open caveats are collected for the reader.

\subsection{Reconciled (code is authoritative)}
\begin{description}[leftmargin=0pt,itemsep=3pt]
  \item[E-step mean update.] The ``corrected'' derivation gives
    $\vm\leftarrow\vm+\Ktil(\Ktil+\GE)^{-1}(\gE-\vm)$; the code uses the older
    $\vm\leftarrow\Ktil(\Ktil+\GE)^{-1}(\GE\Ktil^{-1}\vm+\gE)$
    (\eqref{eq:estep-m}). Both maps share the fixed point $\vm^{*}=\gE$, so the
    \emph{converged} posterior is identical; only the transient differs
    (\S\ref{subsec:estep}).
  \item[Constraint transforms.] $\sigz=\exp(\text{raw})$ and
    $\Amp=\mathrm{softplus}(\text{raw})$; the handoff note claiming ``softplus for
    sigma\_0'' is wrong (\S\ref{subsec:mstep}).
  \item[Two gradient systems.] The ``compute $\mathrm dK$ once, single
    contraction'' form is the kernel-level autograd backward, \emph{not} the
    M-step, which recomputes $\Cmat,\mathrm dC,\Kf,\mathrm dK$ inside the LBFGS
    closure and runs the full moment/loss chain (\S\ref{subsec:mstep}).
  \item[Production utility.] The shipped acquisition is $\Ustd$
    (\texttt{standard\_utility}); the distribution-aware $\Uda$ is the research
    variant derived in the write-ups (\S\ref{sec:utility}).
  \item[Deprecated helper.] \texttt{lambda\_moments} does not exist in the current
    utility code; the moments come from \texttt{get\_gp\_marginal\_moments} and
    \texttt{get\_gp\_conditional\_moments}.
  \item[Cross-covariance.] The manual $\Sigma_{x,x^{*}}$ is unreliable outside the
    inducing hull; the shipped predict path slices GPyTorch's joint
    \texttt{covariance\_matrix}, which already includes $\ccond$
    (\S\ref{sec:crosscov}, Rem.~\ref{rem:cond-correct}).
  \item[Diffusion specs.] The adopted guided-reverse method (Approach~D) uses the
    large linear-schedule model, not the small cosine model; Approach~A optimizes
    by LBFGS, not SGD; some prose describing sizes/status is stale
    (\S\ref{sec:diff-status}).
\end{description}

\subsection{Open caveats (carried honestly)}
\begin{description}[leftmargin=0pt,itemsep=3pt]
  \item[Checkpoint basis drift.] The checkpoint saves $\mb,\Vb$ but not the
    eigenbasis $\eb$, which is recomputed on load; a different torch version can
    yield a subtly different basis and $\sim5\%$ drift (\S\ref{sec:rank1}).
  \item[Utility accuracy ceiling.] At high firing the standard utility can diverge
    and the marginal entropy collapses through the $\rmax$ truncation; a firing
    guard keeps generation in the accurate regime (\S\ref{subsec:numerics},
    \S\ref{sec:diff-status}).
  \item[Subspace offset.] In the C-eigenspace optimization the non-zero offset
    equal to the mean is not mathematically motivated; it compensates a limiter
    precision issue and remains to be investigated (\S\ref{subsec:subspace}).
  \item[Divergence under free ascent.] Because the arc-cosine kernel is positively
    homogeneous, unconstrained gradient ascent on the image inflates the utility
    without bound (Thm.~\ref{thm:divergence}); this is why a subspace/norm
    constraint or a naturalness prior (Part~IV) is required.
\end{description}

\clearpage
\section{Glossary of experimental terms}\label{app:glossary}
\begin{description}[leftmargin=0pt,itemsep=2pt]
  \item[Retinal ganglion cell (RGC)] the output neuron of the retina whose spikes
    are recorded; the GP models one RGC's tuning at a time.
  \item[Receptive field (RF)] the region of the visual field (here, of the image)
    to which the neuron responds; encoded by the structured prior center $\rfc$
    and width $\bwid$.
  \item[Multi-electrode array (MEA)] the recording device that captures spikes
    from many RGCs simultaneously.
  \item[Spike count] the number of spikes $\spk$ a neuron fires in the response
    window to a displayed image; the observed data.
  \item[Firing rate] the instantaneous rate $\rate=e^{\gain\lat+\latz}$; the
    Poisson intensity generating the spike count.
  \item[Tuning] how the neuron's firing rate depends on the stimulus image; the
    latent function $\lat$ the GP infers.
  \item[Natural image] a grayscale photograph-like stimulus (as opposed to
    synthetic gratings); the stimulus domain the prior is designed for.
  \item[Closed loop] the online cycle: display an image, count spikes, update the
    GP, choose the next image.
  \item[Active learning] choosing the next stimulus to maximize expected
    information gain about the neuron's tuning, rather than sampling at random.
\end{description}
